\section{Basic Setup}
To start, you would need a case directory created in your \codeword{CASES} directory. Once created, you will not have a \codeword{caseSetup} file just yet. To get a file, execute the following line in your case directory,
\begin{lstlisting}
    caseSetup --new
\end{lstlisting}
A fresh \codeword{caseSetup} file will be generated. The inside of the file will look like this:
\begin{lstlisting}
    [DEFAULTS]
    DEFAULT_MODULES = defaultGlobalControl defaultGlobalSimControl defaultBCSetup defaultGlobalDecomposition defaultGlobalMaterial defaultGlobalRefinement defaultPost
    ADDON_MODULES =

    [TITLES]
    CASENAME = casename
    CASEDESCRP = default description
    JOBCODE = TS

    [GEOMETRY]
    GEOM = default.stl,1,8,5,1.1,high
\end{lstlisting}
The \codeword{DEFAULT_MODULES} line controls the settings that are available. If the setting is not visible, it just means that it will use the default setup. Each setting here is explained below:
\begin{itemize}
    \item \codeword{CASENAME} - the case name for your case, it does not do anything at this time, but in the future it could be added to the post processing images.
    \item \codeword{CASEDESCRP} - case description, see \codeword{CASENAME}
    \item \codeword{JOBCODE} - Two letter initial for the job you are running, just helps identify the job that's running on the cluster queue.
    \item \codeword{GEOM} - List of geometries to use in the case. A comma separates each setting for that particular geometry, the next geometry goes on the subsequent line.
    \begin{itemize}
        \item What each column represents: 
        \begin{lstlisting}
            GEOMETRY_FILE_NAME,SCALE,MESH_LEVEL,INFLATION_LAYERS,INFLATION_LAYER_GROWTH_RATE,HIGH OR LOW REYNOLDS WALL MODEL
        \end{lstlisting}
    \end{itemize}
\end{itemize}